\begin{table}[htbp]
\centering
\small
\caption{Escolha da melhor métrica few-shot do Problema A para a transferência (google/gemini-2.5-flash-lite).}
\label{tab:T7_escolha_flashlite}
\begin{tabular}{lrrrrrrrl}
\toprule
Métrica de origem (A) & Fid.\ no A & Métr.\ vs real (A) & cos(Perc., NNLS) & Consist.\ B & $\kappa$ B & Consist.\ C & $\kappa$ C & Eleição \\
\midrule
zero-shot (referência) & 81,3 (78,7–83,3) & 17,6 (15,3–21,3) & 0,998 & 72,7 & 0,263 & 92,3 & 0,846 &  \\
4-shot & 93,8 (63,0–99,3) & 92,6 (83,3–98,0) & 0,926 & 70,9 & 0,257 & 91,6 & 0,830 &  \\
10-shot & 97,5 (87,1–100,0) & 94,7 (88,7–98,0) & 0,837 & 69,6 & 0,249 & 91,4 & 0,828 &  \\
20-shot & 97,9 (95,4–100,0) & 95,4 (92,7–98,0) & 0,815 & 69,8 & 0,249 & 91,3 & 0,826 &  \\
\textbf{40-shot} & \textbf{99,0 (95,5–100,0)} & \textbf{95,3 (92,7–98,0)} & \textbf{0,858} & \textbf{69,6} & \textbf{0,242} & \textbf{91,3} & \textbf{0,826} & \textbf{eleita} \\
\bottomrule
\multicolumn{9}{p{0.95\linewidth}}{\footnotesize Critério de eleição: maior fidelidade média no Problema A (Perceptron, R$^2$, sem colapso) entre as métricas few-shot; o zero-shot entra só como referência e o teto não concorre. Consistência e $\kappa$ em B e C: rótulos zero-shot do LLM contra a métrica de origem (Perceptron, só $\hat{W}$ transferido, centróides recalculados pelos rótulos do LLM). Em negrito: a métrica eleita para a transferência.} \\
\end{tabular}
\end{table}
